%==============================================================================
% 3. OrgMemBench Design  (the core contribution)
%==============================================================================
\section{\bench{} Design}
\label{sec:design}

\subsection{Overview: Goals and Scope}
\label{sec:design:overview}

\bench{} is designed to measure one capability that existing benchmarks do not
reach: an agent's ability to answer everyday organizational questions over a
multi-source, multi-author, multi-year record of work. Concretely, the target
queries are the kind a person routinely routes to an assistant with access to
the company's history---``who owns the Acme relationship, and where was that
established?'', ``what did engineering decide about the migration last quarter,
and has that changed?''
Answering these requires stitching evidence across channels (email, Slack,
meeting transcripts, tickets, documents, CRM entries), across people whose roles
and seniority shift over time, and across time itself---including facts that
were later superseded, corrected, held in unresolved contradiction, or never
written down explicitly as a rule. The benchmark is built around four design
principles that follow from this target: ground truth is projected
deterministically from a known source-of-truth graph rather than read back out
of free text; the temporal model is bi-temporal so that ``what was true then''
can be separated from ``what the organization recorded, and when''; scoring is
rubric-based facet coverage rather than exact match, so that multiple correct
phrasings earn credit while omissions are penalized deterministically; and the
harness is \emph{capability-aware}, reporting a category as not-applicable for a
system whose design cannot in principle answer it (for example, an as-of-date
temporal query posed to a system with no temporal model) rather than silently
scoring it zero.

It is equally important to state what \bench{} does \emph{not} test. It is not a
conversational or personal-memory benchmark (the regime LoCoMo, LongMemEval, and
BEAM cover); the full scope boundaries are enumerated in
\S\ref{app:construction}.

\subsection{The Helix Organization and the Bi-temporal Model}
\label{sec:design:helix}

The seed organization is \textbf{Helix Logistics}, a small-to-mid-size
SaaS/freight-tech company instantiated with a coherent multi-year narrative arc:
founded in 2020 by Maya Patel and Daniel Kim, it ships its API and pivots toward
mid-market, weathers a customer escalation and reorganization, and overhauls its
documentation, across a dated 2020--2026 timeline. A fixed company skeleton pins
the founding year, founders, per-year arcs, documentation-discipline eras,
product lines, and a tooling timeline with explicit replacements (for example,
Salesforce replaces HubSpot and Confluence replaces Notion in successive years);
those replacements are not decoration---they seed the temporal and supersession
structure that the hard question categories then probe. The full company
skeleton is enumerated in \S\ref{app:construction}.

The substrate that makes this possible is a \textbf{bi-temporal}
source-of-truth graph. Every event and fact carries four timestamps:
\texttt{valid\_at} (when it became true in the world), \texttt{invalid\_at}
(when it stopped being true), \texttt{ingested\_at} (when the organization
recorded it), and \texttt{expired\_at} (when that record was retracted or
corrected). Two timelines therefore coexist---\emph{valid time} (the world) and
\emph{transaction time} (the record)---and a question can interrogate either or
both. Facts are \textbf{invalidated, never deleted}, which preserves a no-loss
audit trail: the prior state of any fact remains recoverable alongside its
replacement. When one fact replaces another, a \textbf{supersession link}
connects the old fact to the new one together with an explicit \emph{reason}, so
the change is not merely detectable but explainable. Finally, the model supports
\textbf{\texttt{as\_of} date-travel}: a query can reconstruct what the
organization believed on a specified past date, which is distinct from what is
true now. This structure licenses the C1/C3/C4 categories below.

\subsection{Question Taxonomy}
\label{sec:design:taxonomy}
\label{sec:design:questions}

{\sloppy
Each tier ships a versioned question set spanning six structured categories
(C1--C6; Table~\ref{tab:categories}), each isolating a distinct reasoning pattern
and pinned to the facets a correct answer must recover. A seventh,
\textsc{Emergent}, category (reconstructing a de-facto pattern never written down
as an explicit rule) is deferred to the large tier and is out of scope here. The
small tier poses $11$ questions over a $121$-artifact corpus ($\approx$60K
tokens); the medium tier poses $73$ questions over $443$ artifacts
($\approx$200K tokens)---the unit of scale being the corpus a system must ingest
and retrieve over, not the question count alone. Question counts per category and
tier are given in Appendix~\ref{app:construction}
(Table~\ref{tab:design:counts}).
\par}

\begin{table}[t]
  \centering
  \caption{\textbf{Category definitions.} The six structured categories
    evaluated at the small and medium tiers. A seventh category, emergent
    pattern, is evaluated only at the large tier and above and is out of scope
    here.}
  \label{tab:categories}
  \begin{tabular}{l l p{0.46\textwidth}}
    \toprule
    Code & Name & What it tests \\
    \midrule
    C1 & Supersession        & The current value of a fact plus what it
                               replaced, when, and why. Example: \emph{``What is
                               the data-retention policy now, what did it
                               replace, and what drove the change?''} \\
    C2 & Decision provenance & Who decided, when, what alternatives were
                               weighed, and the deciding rationale. Example:
                               \emph{``Who chose the cloud-native migration and
                               what options were rejected?''} \\
    C3 & Bi-temporal (as-of) & What the organization believed AS OF a past date,
                               distinct from now. Example: \emph{``What would the
                               answer have been on 2022-01-15?''} \\
    C4 & Audit replay        & Reconstruct the knowledge state at a past date and
                               flag what has since changed. \\
    C5 & Justification chain & Enumerate the evidence supporting a conclusion and
                               classify each as direct testimony vs. inference. \\
    C6 & Contradiction       & Detect that two artifacts conflict; report both
                               sides and whether it was resolved. \\
    \bottomrule
  \end{tabular}
\end{table}

The ground-truth and rubric field schemas, the $72$-persona library and
document-genre taxonomy, and the artifact counts behind these question sets are
detailed in Appendix~\ref{app:construction}.

\subsection{Synthetic Generation}
\label{sec:design:why-synthetic}
\label{sec:generation}
\label{sec:design:generation}

The data is synthetic by deliberate design: a freshly generated fictional
organization cannot appear in any system's pretraining data or retrieval cache, so
a high score cannot be explained by memorization---a stronger guarantee than
post-hoc $n$-gram decontamination, which can only estimate leakage that has
already occurred (the further reasons of control and privacy are developed in
Appendix~\ref{app:construction}). Ground-truth labels are deterministic
projections of the source-of-truth graph rather than free-text recall, which
sharply limits label ambiguity but does not remove the single-model circularity
inherent in one model writing the artifacts, phrasing the questions, and running
the self-consistency check; following \citet{gebru2021datasheets} we document
this rather than obscure it (in a full datasheet, Appendix~\ref{app:datasheet}),
and a cross-model validation pass is planned. We accompany the release with
Croissant metadata~\citep{croissant2024} so the dataset is machine-discoverable
and ML-ready. The \texttt{helix-corpus} generation tooling (a single local
\texttt{gpt-oss:20b} model served via Ollama), the full Stage~A--F pipeline, and
the logged, reversible three-wave manual audit that reconciled shipped files with
the graph without altering the answer key are detailed in
Appendix~\ref{app:construction}; the audit log and its six defect classes in
Appendix~\ref{app:annotation}.
